
Execution feedback enables code repair models to achieve higher bug fix rates than global rewrites, yet most refinement methods treat feedback as a black-box signal. When a language model generates buggy code, detailed error traces---line numbers, error type, expected versus actual values---provide localization information that differs from the approximate signals AI critics can offer. Understanding this difference has practical implications for practitioners designing iterative refinement pipelines, who face a choice between execution sandboxes and LLM-based critics with different infrastructure cost, latency, and repair success characteristics.

Prior work has demonstrated improvements from both execution-based approaches (Chen et al., 2023; Le et al., 2022) and AI-based approaches \citep{madaan2023self}. However, comparisons between these methods confound the feedback signal with differences in model architecture, prompting strategy, and refinement loop design. This study addresses this gap by fixing the method (prompt-based iterative refinement with CodeLlama-7B-Instruct) and varying only the feedback signal across four conditions: execution-detailed, execution-binary, AI-critic, and random baseline.

The central hypothesis is that execution feedback's advantage arises from information structure rather than information quantity---specifically, error traces contain counterfactual localization (''if X were different, Y would not have failed'') that constrains the model's hypothesis space for edits. This study tests this hypothesis through a three-step causal chain: (1) execution traces contain extractable counterfactual information, (2) this information enables targeted rather than global edits, and (3) targeted edits achieve higher fix rates.

The experiments validate five of six sub-hypotheses. The complexity interaction hypothesis (H-C1) is refuted---execution advantage does not increase with task complexity, contrary to initial predictions.
